\chapter{Metodología y arquitectura del sistema}
\label{ch:metodologia}
% Mapa de fuentes por sección: MAPA_FUENTES.md, capítulo 4.
\section{Del desarrollo exploratorio al código reutilizable}
% README.md; notebooks/01,03,07,10,11,12; src/renewables_permitting/pipeline.py.
\pendiente{Explicar la exploración con notebooks y el traslado de lógica
estable al paquete reutilizable: modularidad y separación de responsabilidades.
No afirmar que todos los notebooks históricos se migraron por completo.}

\section{Organización en capas Bronze, Silver y Gold}
% docs/USER_GUIDE.md §§2--3; docs/architecture/silver_extraction_contract.md.
\pendiente{Definir primero datos próximos a la fuente, tablas estructuradas
y datos consolidados para consulta; después introducir Silver y Gold.
Puede introducirse la analogía materia prima, información estructurada y
producto de consulta, seguida de la definición exacta de cada capa.
Explicar ETL (extraer, transformar y cargar datos) solo si se utiliza el término;
no confundir esa extracción general con la extracción de entidades por IA.
Situar extracción, 13 tablas Silver, derivación territorial y cuatro tablas
Gold. Reservar una figura global con una rama experimental separada.}

\section{Extracción estructurada mediante Gemini}
% extraction/config.py, documents.py, models.py, instructions.py; resultados §1.
\pendiente{Explicar por qué se utiliza un LLM, qué recibe y qué devuelve.
Describir Gemini 2.5 Flash como parte del sistema evaluado; separar selección
del texto, instrucciones, evidencia y límites de inferencia. TODO-CITA:
documentación del modelo y extracción estructurada.}

\section{Contratos, validación y canonicalización}
% architecture/silver_extraction_contract.md; extraction/models.py,
% flat_contract.py, canonicalization.py, validation.py.
\pendiente{Definir un contrato como la estructura y los valores admisibles;
Pydantic es la biblioteca que permite expresarlo y comprobarlo. Introducir
categorías permitidas, nulos, relaciones y evidencia. Distinguir generación
por IA de reglas deterministas y validez estructural de corrección semántica.
Reservar la figura de extracción y validación. TODO-CITA: Pydantic.}

\section{Agrupación de menciones en proyectos}
% FINAL_CORPUS_W14_DOWNSTREAM_MATERIALIZATION.md §5; project_grouping.py;
% tests/test_project_grouping.py; FINAL_W14_ADMIN_ACTION_TEMPORAL_AUDIT.md.
\pendiente{Definir activo de generación, mención e identidad consolidada;
introducir después la clave project\_match\_key. Explicar nombres y alias
observados, tecnología y anclaje territorial; solo la generación crea raíces.
Reservar un esquema de menciones documentales que convergen en un proyecto.
Reutilizar Don Rodrigo II sin confundir agrupación longitudinal con el
emparejamiento de predicciones y referencia en la evaluación.}

\section{Resolución territorial}
% USER_GUIDE.md; FINAL_CORPUS_W14_DOWNSTREAM_MATERIALIZATION.md §§3--4;
% ine_reference.py, location_resolution.py, project_locations.py.
\pendiente{Explicar la correspondencia entre topónimos y códigos oficiales
INE de municipio, provincia y comunidad o ciudad autónoma. Conservar
resoluciones parciales y ambigüedades; territorio documental no equivale a
ubicación exacta de la planta. TODO-CITA: referencia INE realmente utilizada.}

\section{Revisión humana y antecedentes históricos}
\subsection{Decisiones y correcciones versionadas}
% architecture/extraction_quality_review.md; FINAL_W14_ADMIN_ACTION_CORRECTIONS.md.
\pendiente{Definir cola de revisión, decisiones de la autora y correcciones
versionadas. Explicar conservación del original, aplicación antes de Silver,
trazabilidad y regeneración sin editar tablas derivadas.}
\subsection{Salvaguarda de antecedentes históricos}
% architecture/extraction_quality_review.md §Safeguard;
% extraction/historical_antecedents.py; tests/extraction/test_historical_antecedents.py.
\pendiente{Introducir P0 como nombre de la salvaguarda: requiere una señal
temporal o estructural y otra contextual independiente. Abre revisión sin
eliminar actuaciones. Separar detección, decisión humana y corrección;
no confundirla con la recuperación retrospectiva del corpus.}

\section{Reproducibilidad y proceso de desarrollo}
\subsection{Interfaz de línea de comandos y pruebas}
% README.md; USER_GUIDE.md §§5,14; pipeline.py; .github/workflows/tests.yml.
\pendiente{Definir CLI como una interfaz de instrucciones por terminal y
pytest como herramienta de pruebas automatizadas. Explicar repetición de
transformaciones, contratos, regresiones y casos límite. No enumerar todos
los comandos ni usar el número de tests como métrica científica.}
\subsection{Desarrollo asistido por IA}
\label{sec:ia-desarrollo}
% Declaración de uso aportada por la autora en la petición de esta fase;
% AGENTS.md; FINAL_W14_ADMIN_ACTION_CORRECTIONS.md §1; resultados §§1,11.
\pendiente{Describir, según la declaración de la autora, ChatGPT para análisis,
planificación, diseño metodológico, revisión y refinamiento de prompts;
Codex desde VSCode para inspección, implementación acotada, pruebas,
refactorización, documentación y LaTeX. Definir primero el archivo de
instrucciones persistentes que aporta reglas y contexto al agente; después
nombrar AGENTS.md y explicar su uso junto a prompts con alcance y restricciones;
desarrollo incremental, revisión humana, tests y commits/push manuales de la
autora, con separación de ramas. Contrastar la cronología de las versiones
fijadas de sistema, referencia humana y evaluador antes de las predicciones.
No hubo ajuste retrospectivo del experimento. Distinguir estos asistentes de
Gemini como objeto evaluado. La autora decide objetivos y método, anota,
revisa, interpreta y asume la responsabilidad; Codex no es autor del TFM.
La declaración de uso debe ser revisada por ella y no se deduce de Git.}
